% !TeX root = ../main.tex
% 系统测试设计与计划

\chapter{系统测试设计与计划}
软件测试作为系统开发生命周期中的关键环节，对于保障系统功能的完整性和性能的可靠性具有至关重要的作用。通过科学合理的测试设计与执行，不仅能够有效发现潜在缺陷，提升系统的稳定性与健壮性，还能为后续的维护与优化提供有力依据。

本章将围绕测试目标与设计思路展开，重点描述功能测试和性能测试方面的计划与方法。鉴于涉及的测试数据与结果具有保密要求，本章主要阐述系统测试的设计与总体规划，不予披露具体测试过程与测试结果。

\section{系统开发环境}
系统开发环境与主要工具选型如表~\ref{tab:6.1}所示。开发环境包括操作系统、开发工具、编程语言、设计工具及硬件平台等多个方面，确保系统从设计、实现到部署的全流程支持。

\begin{table}[ht]
\centering
\caption{系统开发环境}
\label{tab:6.1}
    \begin{tabular}{ll}
        \toprule
        \textbf{类别} & \textbf{具体工具/平台} \\ 
        \midrule
        操作系统 & Linux  \\
        开发工具 & Visual Studio Code, WebStorm \\ 
        编程语言 & C, Python, Vue.js  \\ 
        设计工具 & PlantUML, Draw.io \\
        硬件平台 & PC, 服务器, 嵌入式单板 \\
        \bottomrule
    \end{tabular}
\end{table}

\begin{itemize}
\item \textbf{操作系统}：选用 Linux 作为开发与部署环境，为系统提供稳定的多任务调度能力、丰富的工具链支持以及对嵌入式交叉编译环境的兼容性。Linux 环境有利于后续自动化构建、测试及跨平台运行。
\item \textbf{开发工具}：采用 Visual Studio Code 与 WebStorm。前者提供对 C 语言和 Python 的高效代码编辑与调试功能，支持丰富的插件生态和版本控制集成；后者针对前端 Vue.js 提供语法高亮、代码补全、调试以及构建工具集成，提高前端开发效率与代码质量。
\item \textbf{编程语言}：系统主要使用 C、Python 与 Vue.js。C 语言用于覆盖率采集组件的高性能实现，Python 用于后端服务逻辑与数据解析处理，Vue.js 用于前端界面可视化呈现，实现模块化组件化开发，以提高维护性和扩展性。
\item \textbf{设计工具}：选用 PlantUML 与 Draw.io，用于系统架构设计、类图、流程图等的绘制。通过可视化设计工具，能够清晰描述系统逻辑与模块交互关系，降低设计误差，提高开发效率。
\item \textbf{硬件平台}：包括 PC、服务器及嵌入式单板。PC 用于本地开发与测试，服务器用于系统部署与服务支撑，嵌入式单板作为目标平台，用于覆盖率数据采集。多平台组合保证系统开发、调试与验证环节的完整性。
\end{itemize}

\section{系统测试}
\subsection{系统测试概览}
系统测试计划概览如表~\ref{tab:6.2}所示。概括了测试环境、硬件配置、测试数据、测试方法、测试目标以及测试范围等关键要素。

\begin{table}[ht]
\centering
\caption{系统测试概览}
\label{tab:6.2}
    \begin{tabular}{lp{10cm}}
        \toprule
        \textbf{类别} & \textbf{内容说明} \\ 
        \midrule
        测试环境 & PC（Linux）、嵌入式单板 \\
        硬件配置 & PC: x86\_64, 8GB 内存；单板: 多种不同架构的硬件平台 \\ 
        测试数据 & 生产环境中的真实覆盖率数据 \\ 
        测试方法 & 单元测试、集成测试、系统测试、性能测试 \\ 
        测试目标 & 验证系统是否满足功能/非功能需求 \\ 
        测试范围 & \begin{itemize}[leftmargin=*]
            \item 覆盖率采集模块：覆盖率采集的完整性以及跨操作系统共享内存通信功能
            \item 细粒度分析模块：各领域原始覆盖率数据的解析能力，以及模块覆盖率、增量覆盖率计算的准确性
            \item 覆盖率可视化模块：可视化界面及其交互逻辑的可用性
            \item SmartCI集成：任务触发、执行及结果回传的正确性与可靠性
        \end{itemize} \\
        \bottomrule
    \end{tabular}
\end{table}

\begin{itemize}
    \item \textbf{测试环境}：系统在 PC（Linux 平台）及多种架构的嵌入式单板上进行验证，确保测试覆盖了开发与部署所涉及的主要硬件环境。
    \item \textbf{硬件配置}：PC 平台采用 x86\_64 架构并配备 8GB 内存，嵌入式单板包括多种不同架构，以确保系统在多样化硬件条件下的稳定性与兼容性。
    \item \textbf{测试数据}：采用生产环境中的真实覆盖率数据，确保测试结果能够真实反映系统在实际应用中的表现。
    \item \textbf{测试方法}：采用多层次测试策略，包括单元测试、集成测试、系统测试以及性能测试，全面验证各模块的功能完整性、接口交互及性能表现。
    \item \textbf{测试目标}：评估系统是否满足功能需求和非功能需求。
    \item \textbf{测试范围}：按照功能模块划分，测试各个功能模块的功能需求和非功能需求是否满足。
\end{itemize}

\section{系统测试设计}
测试设计以需求分析为导向，确保覆盖率数据处理、可视化展示及自动化执行等关键功能的完整性与可靠性。通过测试用例设计的方式，详细描述系统测试的设计与计划。

鉴于保密策略，本节无法写出完整的可执行测试用例，包括输入的具体文件、期望数值等。因此进行抽象化用例设计，不涉及实际数据，仅描述测试对象、输入类型、操作和预期结果。

\subsection{覆盖率采集模块}
覆盖率采集模块的测试用例设计如表~\ref{tab:6.3}所示。主要围绕数据采集的功能完整性和跨操作系统的数据传输可靠性展开。

\begin{table}[ht]
\centering
\caption{系统测试用例设计-覆盖率采集模块}
\label{tab:6.3}
\begin{tabular}{p{2cm}p{2.5cm}p{4cm}p{4cm}}
\toprule
\textbf{测试对象} & \textbf{输入} & \textbf{操作步骤} & \textbf{预期结果} \\
\midrule
覆盖率采集 & 插桩代码 & 调用采集接口，收集插桩统计数据 & 返回覆盖率数据完整，文件格式正确 \\
覆盖率采集 & 模拟数据面数据 & 模拟数据面采集数据并通过共享内存传输 & 数据传输正确，汇总结果一致 \\
覆盖率采集 & 冒烟测试触发信号 & 执行自动采集流程 & 数据采集完整，无丢失 \\
\bottomrule
\end{tabular}
\end{table}

首先，重点验证采集接口及插桩数据收集函数的功能正确性。测试通过模拟插桩数据输入，调用系统提供的采集接口，并收集生成的插桩数据，以检验返回的数据是否完整，生成的覆盖率文件是否符合预期设计规范。通过此类测试，可以确保采集模块的基本功能符合开发需求，避免因接口异常导致的数据丢失或格式错误问题。

其次，测试重点转向跨操作系统的共享内存通信功能的正确性。测试模拟无文件系统的数据面的覆盖率数据，通过共享内存机制进行传输，随后在主机端进行汇总和验证。通过对比传输前后汇总的覆盖率数据文件，判断数据传输的完整性和一致性，确保系统能够在异构操作系统架构下稳定运行，并能够正确汇总异构操作系统采集的数据。

最后，覆盖率采集模块在真实冒烟测试触发场景下进行验证。测试通过冒烟测试触发信号启动自动采集流程，模拟实际运行环境下的采集任务，并对采集结果进行完整性检查。该测试确保系统在实际部署场景中能够按照预定流程自动采集覆盖率数据，保证数据的完整性和可靠性，为后续分析模块提供准确、完整的原始数据基础。

\subsection{细粒度分析模块}
细粒度分析模块的测试用例设计如表~\ref{tab:6.4}所示，主要针对覆盖率数据的解析精度、模块覆盖率及增量覆盖率计算的准确性，以及告警信息的正确生成。

\begin{table}[ht]
\centering
\caption{系统测试用例设计-细粒度分析模块}
\label{tab:6.4}
\begin{tabular}{p{2cm}p{2.5cm}p{4cm}p{4cm}}
\toprule
\textbf{测试对象} & \textbf{输入} & \textbf{操作步骤} & \textbf{预期结果} \\
\midrule
细粒度分析 & 各领域覆盖率原始文件 & 全量覆盖率解析 & 覆盖率解析正确，覆盖率报告无异常 \\
细粒度分析 & 全量覆盖率信息 & 调用模块、增量覆盖率解析方法 & 模块、增量覆盖率解析正确，报告无异常 \\
细粒度分析 & 解析后的覆盖率信息 & 调用告警计算方法 & 告警信息输出正确 \\
细粒度分析 & 各领域覆盖率原始文件 & 执行完整解析流程 & 解析正确，覆盖率数据一致 \\
\bottomrule
\end{tabular}
\end{table}

首先，针对各领域的覆盖率原始文件的解析能力进行验证，确保全量覆盖率解析功能能够正确提取数据，并生成规范化的覆盖率报告。测试通过调用解析函数处理真实的覆盖率数据文件，检查解析结果与预期是否一致，同时验证生成报告是否存在异常或缺失内容。该步骤的目标是保证覆盖率解析算法能够正确识别每个领域和模块的覆盖情况，为后续分析提供准确的数据基础。

完成全量覆盖率解析测试后，重点验证模块覆盖率和增量覆盖率的计算方法。通过调用对应解析方法，处理不同领域、不同基准节点的全量覆盖率数据，检查解析结果是否准确计算出模块覆盖率以及增量覆盖率。

在告警信息生成的测试中，使用解析后的覆盖率信息调用告警计算方法，验证告警输出的正确性。该测试目标是确保系统能够及时、准确地发现模块配置错误或函数未覆盖的情况，支持质量保障和问题定位。

最后，执行完整解析流程，从各领域原始覆盖率文件开始，经过覆盖率解析、增量计算及告警生成，最终生成完整的覆盖率数据与报告。测试目标包括验证数据处理流程的连贯性、各环节解析结果的一致性，以及覆盖率数据和报告的完整性，确保细粒度分析模块在实际运行环境下能够提供准确、可靠的分析结果，为可视化展示和决策提供坚实的数据支撑。

\subsection{覆盖率可视化模块}
覆盖率可视化模块的测试用例设计如表~\ref{tab:6.5}所示，主要验证前端界面的完整性、数据交互的流畅性以及用户操作的有效性。

\begin{table}[ht]
\centering
\caption{系统测试用例设计-覆盖率可视化模块}
\label{tab:6.5}
\begin{tabular}{p{2.5cm}p{2.5cm}p{3.5cm}p{4cm}}
\toprule
\textbf{测试对象} & \textbf{输入} & \textbf{操作步骤} & \textbf{预期结果} \\
\midrule
覆盖率可视化 & 覆盖率数据 & 加载各组件并绑定数据 & 表格、柱状图、告警列表正确显示 \\
覆盖率可视化 & 覆盖率数据 & 切换维度，刷新数据 & 数据刷新正常，筛选功能有效 \\
覆盖率可视化  & 覆盖率数据 & 全局操作界面、筛选、交互 & 界面完整，交互逻辑正确 \\
\bottomrule
\end{tabular}
\end{table}

首先，进行加载各组件并绑定数据的测试，通过真实的覆盖率数据输入，验证表格、柱状图和告警列表等组件的显示效果。测试目标是确保这些组件在渲染过程中能够正确地显示覆盖率数据，确保各数据维度得到清晰、准确的展示，避免数据展示的遗漏或错误。这一部分的重点是确保界面的基础功能得到可靠执行，并验证组件间的数据传递和更新是否无误。

接下来的测试将验证切换维度和数据刷新功能的正常运行。通过模拟用户切换不同数据维度（如时间、单板、领域等）进行筛选和展示，检查数据的刷新情况。测试目标是确保系统能够根据用户操作及时、准确地刷新展示的数据，并验证筛选功能在不同数据维度下的有效性。重点关注各个组件的响应是否出现异步，刷新过程中是否出现性能瓶颈，确保界面更新不受影响，且数据始终与用户选择的维度保持一致。

最后，进行全局用户界面的完整性及交互逻辑验证。测试通过模拟不同的用户操作场景，确保系统支持用户执行各种交互操作，如筛选、查看报告、选择维度等，并验证这些操作是否能够准确反馈至系统，保持数据一致性。测试目标是确保前端界面不仅在功能上完整，而且用户交互体验流畅、逻辑正确，符合设计要求。重点检查界面组件与底层数据交互的正确性，确保用户通过简单、直观的操作能够获得准确的覆盖率分析结果。

\subsection{数据存储模块}
数据存储模块的测试用例设计如表~\ref{tab:6.6}所示。主要针对数据库的插入与更新，查询操作的正确性，静态文件系统中的报告存储及访问功能，以及过期数据的清理功能。

\begin{table}[ht]
\centering
\caption{系统测试用例设计-数据存储模块}
\label{tab:6.6}
\begin{tabular}{p{2.5cm}p{2.5cm}p{3.5cm}p{4cm}}
\toprule
\textbf{测试对象} & \textbf{输入} & \textbf{操作步骤} & \textbf{预期结果} \\
\midrule
数据存储模块 & 覆盖率数据 & 向数据库插入/更新数据 & 数据库成功响应，在数据表中能找到新的数据 \\
数据存储模块 & 数据库数据 & 从数据库查询数据 & 数据库成功响应，能查询到覆盖率记录 \\
数据存储模块 & 覆盖率报告 & 向静态文件系统存储报告 & 静态文件系统中存在文件，并能通过URL访问 \\
数据存储模块 & 过期覆盖率数据和报告 & 调用过期数据清理方法 & 过期数据和报告被删除 \\
\bottomrule
\end{tabular}
\end{table}

首先，进行向数据库插入或更新数据的测试，主要验证系统是否能够正确响应插入请求，并确保数据被正确存储到数据库中。在测试过程中，输入一组覆盖率数据，系统应当能够成功处理该数据并将其存储至数据库，数据库中能够正确反映输入数据，并确保各字段值与原始数据一致，以此验证数据存储的准确性和数据库的响应能力。

其次，数据库查询操作是验证系统能否高效并准确地返回已存储数据的关键。通过模拟从数据库中查询覆盖率数据的操作，检查系统是否能够成功返回正确的结果。测试应确保从数据库查询到的记录与存储时的数据完全一致，从而验证数据库查询功能的完整性与准确性。

对于覆盖率报告的存储与访问，测试的重点在于静态文件系统的存储功能是否稳定可靠。通过将覆盖率报告存储到静态文件系统中，并尝试构造静态资源 URL 访问该报告，来验证文件存储与读取是否成功。预期结果应当是文件能够成功存储，并且用户通过指定的 URL 地址能正确地访问文件，这确保了文件系统的正常运作以及数据的可访问性。

最后，测试过期数据清理是否正常运行，验证系统是否能够准确识别并删除过期的数据与报告。在测试时，手动写入过期的覆盖率文件和覆盖率数据，调用清理方法进行清理。系统应能够正确地判断数据的过期状态，删除已经过时的文件或记录，而不影响有效数据的存在。这一功能的测试确保了系统的资源管理能力，并避免了无用数据的积累。

\subsection{SmartCI 集成}
SmartCI 集成的测试用例设计如表~\ref{tab:6.7}所示。主要针对系统在自动化任务执行中的正确性与效率，特别是任务的执行与数据回传的准确性。

\begin{table}[ht]
\centering
\caption{系统测试用例设计表}
\label{tab:6.7}
\begin{tabular}{p{2.5cm}p{2.5cm}p{3.5cm}p{4cm}}
\toprule
\textbf{测试对象} & \textbf{输入} & \textbf{操作步骤} & \textbf{预期结果} \\
\midrule
SmartCI 集成  & 冒烟测试触发信号 & 执行流水线任务 & 覆盖率采集/分析/可视化结果正确回传 \\
SmartCI 集成  & 冒烟测试触发信号 & 执行整个流水线 & 覆盖率任务自动执行完成，数据正确 \\
\bottomrule
\end{tabular}
\end{table}

首先，测试的重点是单个任务的执行。具体来说，构造单个 task 组成的 project，系统接收到冒烟测试触发信号后，执行单个任务单元，如覆盖率采集、分析或可视化任务。此时，测试的目的是验证任务单元本身的功能是否能够正确执行，并确保各任务模块在独立运行时能够正常返回预期结果。在执行该任务时，系统应能够成功触发任务并完成相应的效果，最终确保所有数据能正确回传并呈现给用户。

在测试整个流水线的执行时，系统从冒烟测试触发信号开始，启动完整的自动化流水线，依次执行所有任务单元。测试目标是确保各任务单元能够按预定顺序正确执行，且任务间的数据流与依赖关系能够无误传递，确保整个流程的正确性和稳定性。重点验证任务的协调性，确保覆盖率采集、数据分析与可视化展示的无缝衔接。预期结果是整个流水线顺利执行，任务间的数据传递与执行过程无误，最终准确生成覆盖率数据和报告，自动化执行符合预期。

\subsection{性能测试}
性能测试的测试用例设计如表~\ref{tab:6.8}所示。主要测试整个系统的性能，验证系统在大规模数据处理和高负载条件下稳定性和响应性。

\begin{table}[ht]
\centering
\caption{系统测试用例设计-性能测试}
\label{tab:6.8}
\begin{tabular}{p{2.5cm}p{2.5cm}p{3.5cm}p{4cm}}
\toprule
\textbf{测试对象} & \textbf{输入} & \textbf{操作步骤} & \textbf{预期结果} \\
\midrule
整体系统性能 & 大规模覆盖率数据 & 测量采集/解析/渲染时间 & 响应时间在可接受范围内，性能符合预期 \\
数据库性能 & 大规模覆盖率数据 & 测试数据库读写操作的响应时间 & 数据库能在合理时间内完成大规模数据的读取与写入 \\
静态文件系统性能 & 大规模报告存储 & 测试报告存储与访问的响应时间 & 静态文件系统能高效存储和快速访问大量报告 \\
前端系统性能 & 大规模覆盖率数据 & 测试界面加载和数据渲染的响应时间 & 界面渲染流畅，数据加载与刷新及时响应 \\
\bottomrule
\end{tabular}
\end{table}

首先，进行整体系统性能测试，关注大规模覆盖率数据处理的效率，特别是数据采集、解析与渲染的响应时间。目标是确保系统在高负载下能在可接受的时间内完成整个流程，避免性能瓶颈和延迟，确保系统稳定性和流畅性。

其次，进行数据库性能测试，关注大规模数据的读写响应时间。通过模拟大量覆盖率数据的插入、更新和查询操作，验证数据库在高并发条件下的处理能力，确保数据库能够快速完成大规模数据操作，避免性能下降或事务失败。

随后，进行静态文件系统性能测试，验证在大规模报告存储与访问场景下，静态文件系统的效率。测试目标是确保文件系统能高效存储并快速响应大量报告的访问，保证用户体验和系统效率不受影响。

最后，进行前端系统性能测试，关注用户界面在大规模数据加载和渲染时的响应时间。目标是确保系统在合理时间内加载数据，且界面渲染流畅，交互响应及时，以保证优质的用户体验。

\section{本章小结}
本章详细阐述了系统的测试设计与计划，涵盖了功能需求和非功能需求的全面测试。测试计划通过对覆盖率采集、细粒度分析、可视化模块、数据存储、SmartCI 集成等多个核心模块的测试用例设计，确保系统在不同场景下的功能完整性和稳定性。同时，针对系统性能、数据库读写、文件系统存储与访问等关键性能指标进行了细致的性能测试计划，确保系统在大规模数据处理时的高效性与响应速度。